\section{The Conjecture Capsule (Internal Bottlenecks)}
\label{sec:conjecture-capsule}
%=============================================================================
% This section centralizes the manuscript's internal conjectural inputs ("B")
% that are repeatedly invoked across the roadmap.  Each item is cross-referenced
% to the Open Problem Ledger (Appendix~\ref{app:open-problem-ledger}).
%=============================================================================

\begin{tcolorbox}[colback=blue!3!white, colframe=blue!60!black, title=\textbf{How to read this section}]
The goal of this manuscript is to separate what is (i) proved here, (ii) available in the literature with precise hypotheses, and (iii) currently \emph{open}.

This section records the minimal conjectural inputs repeatedly used in later ``if--then'' statements. Each conjecture below is labeled and mapped to one or more items in the Open Problem Ledger (Appendix~\ref{app:open-problem-ledger}).
\end{tcolorbox}

\subsection{B.1. Positivity of string tension beyond strong coupling}

\begin{conjecture}[All-coupling string tension positivity]
\label{conj:sigma-all-coupling}
In four-dimensional $SU(N)$ lattice Yang--Mills theory (Wilson action), the infinite-volume string tension exists and is strictly positive for every $\beta>0$:
\[
\sigma(\beta) > 0, \qquad \forall\,\beta>0.
\]
\end{conjecture}

\begin{remark}[Status and relation to the ledger]
At present, positivity of $\sigma(\beta)$ is rigorously established only in a strong-coupling regime (via cluster expansion). Extending $\sigma(\beta)>0$ to all $\beta>0$ is treated as an \textbf{open input} in this manuscript; see the ledger items responsible for ``confinement at all couplings''.
\end{remark}

\subsection{B.2. Uniform spectral gap mechanism (finite-volume to infinite-volume)}

\begin{conjecture}[Uniform-in-volume gap]
\label{conj:uniform-gap}
Fix $\beta>0$. There exists $\Delta_*(\beta,N)>0$ such that the finite-volume spectral gap of the transfer matrix satisfies
\[
\Delta_L(\beta) \ge \Delta_*(\beta,N) \quad \text{for all volumes } L.
\]
\end{conjecture}

\begin{remark}[Interpretation]
This is the standard ``no gap closing with $L$'' requirement needed to pass from finite-volume positivity $\Delta_L(\beta)>0$ to a nontrivial infinite-volume mass gap bound.
\end{remark}

\subsection{B.3. Absence of bulk/thermal transitions along the roadmap}

\begin{conjecture}[No bulk transition along the interpolation/RG path]
\label{conj:no-phase-transition}
Along the parameter path(s) used to connect the strong-coupling regime to the weak-coupling/continuum regime (e.g. an RG trajectory, interpolation in the action, or mixed-coupling deformations), there is no phase transition that invalidates the intended continuation of confinement or functional inequalities.
\end{conjecture}

\begin{remark}
This conjecture is a placeholder for a family of ``no first-order transition / no deconfinement'' assumptions frequently invoked in informal arguments; it is not proved in this manuscript.
\end{remark}

\subsection{B.4. Mass gap versus string tension: the key quantitative bridge}

Conjecture~\ref{conj:gap-sigma} (Section~\ref{sec:spectral-lower-bound}) is the canonical statement of the bridge between string tension and spectral gap.

\begin{conjecture}[Mass gap--string tension relation (capsule form)]
\label{conj:gap-sigma-capsule}
Assume $\sigma(\beta)>0$ at some $\beta$. Then there exists $c_N>0$ (depending only on $N$) such that
\[
\Delta(\beta) \ge c_N\,\sqrt{\sigma(\beta)}.
\]
\end{conjecture}

\begin{remark}
This capsule duplicates Conjecture~\ref{conj:gap-sigma} only to highlight its role as a bottleneck conjecture. The discussion in Section~\ref{sec:spectral-lower-bound} should be treated as phenomenological motivation, not a proof.
\end{remark}

\subsection{B.5. Continuum limit existence (constructive control)}

Conjecture~\ref{conj:continuum-limit} (Section~\ref{sec:continuum-program}) is the canonical continuum existence statement.

\begin{remark}[OS reconstruction as an explicit proof obligation]
Even when subsequential limits exist, establishing the Osterwalder--Schrader axioms (in particular reflection positivity and the growth/regularity hypotheses needed for reconstruction) is a distinct, nontrivial step; see Appendix~\ref{app:open-problem-ledger}.
\end{remark}

\subsection{B.6. Mapping to the Open Problem Ledger}

For convenience, we record a non-exhaustive mapping between the capsule conjectures and the ledger items (Appendix~\ref{app:open-problem-ledger}):
\begin{itemize}
    \item Conjecture~\ref{conj:sigma-all-coupling} $\leadsto$ ledger items on all-coupling confinement/string tension.
    \item Conjecture~\ref{conj:uniform-gap} $\leadsto$ ledger item(s) on uniform functional inequalities (e.g. uniform LSI / spectral permanence).
    \item Conjecture~\ref{conj:no-phase-transition} $\leadsto$ ledger item(s) on controlling interpolation paths and excluding bulk transitions.
    \item Conjecture~\ref{conj:gap-sigma-capsule} (equivalently Conjecture~\ref{conj:gap-sigma}) $\leadsto$ ledger item(s) providing a quantitative mass-gap lower bound from confinement.
    \item Conjecture~\ref{conj:continuum-limit} $\leadsto$ ledger item(s) on RG control and OS reconstruction.
\end{itemize}

\subsection{B.7. Mathematical Frameworks (Proposed Solutions)}

We introduce specific analysis frameworks designed to resolve these open problems in Appendix~\ref{app:hard_analysis_frameworks}:
\begin{itemize}
    \item \textbf{Framework A (Dirichlet Forms):} A strategy for OP7 (Continuum Limit) based on Mosco convergence.
    \item \textbf{Framework B (Multiscale Coercivity):} A strategy for OP1 (Uniform LSI) based on renormalization of Log-Sobolev constants.
    \item \textbf{Framework C (Variational Operators):} A strategy for OP5 (Gap vs String Tension) using rigorous transfer matrix spectral analysis.
\end{itemize}
